% -----------------------------------------------------------
\section{Belief}
% -----------------------------------------------------------
	\label{BELIEF}
	
% % ----------------------
% \subsection{See also}
% % ----------------------

% % ----------------------
% \subsection{1828 American Dictionary of the English Language} % *1828
% % ----------------------
	% \label{BELIEF-1828}

	% \begin{compactitem}
		% \item
	% \end{compactitem}

% % ----------------------
% \subsection{Oxford English Dictionary} % *OED
% % ----------------------
	% \label{BELIEF-OED}

	% \begin{compactitem}
		% \item[]
	% \end{compactitem}

% % ----------------------
% \subsection{Scriptural Usage} % *SCRIPTURES
% % ----------------------
	% \label{BELIEF-SCRIPTURES}

	% % -----
	% \subsubsection{Old Testament} % *OT
	% % -----
		% \label{BELIEF-OT}
	
		% % --
		% \paragraph{KJV}
		% % --
			% \label{BELIEF-OT-KJV}
	
			% \begin{tabular}{ll}
				% Total occurrences in the OT & 
			% \end{tabular}
			
			% \begin{compactdesc}
				% \item[]
			% \end{compactdesc}
			
		% % --
		% \paragraph{Hebrew Bible}
		% % --
			% \label{BELIEF-HEBREW}
		
			% %
			% \subparagraph{\texorpdfstring{\he{sample word}}{}}
			% %
				% \label{PAGA} % whatever is in step bible without periods
			
				% \begin{compactdesc}
					% \item[HALOT , PDF ] \textbf{}
						% \begin{compactitem}
							% \item[1]
						% \end{compactitem}
					% \item[BDB , PDF ] \textbf{}
						% \begin{compactitem}
							% \item[1]
						% \end{compactitem}
					% \item[DCH PDF ] \textbf{}
						% \begin{compactitem}
							% \item[1]
						% \end{compactitem}
				% \end{compactdesc}
				
				% \begin{compactdesc}
					% \item[Some OT uses] \textbf{}
						% \begin{compactdesc}
							% \item[]
						% \end{compactdesc}
				% \end{compactdesc}
			
		% % --
		% \paragraph{Septuagint}
		% % --
			% \label{BELIEF-LXX}
		
			% %
			% \subparagraph{\texorpdfstring{\gr{sample word}}{}}
			% %
		
	% % -----
	% \subsubsection{New Testament} % *NT
	% % -----
		% \label{BELIEF-NT}
	
		% % --
		% \paragraph{KJV}
		% % --
			% \label{BELIEF-NT-KJV}
			
	
			% \begin{tabular}{ll}
				% Total occurrences in the NT & 
			% \end{tabular}
			
			% \begin{compactdesc}
				% \item[]
			% \end{compactdesc}
			
		% % --
		% \paragraph{Greek New Testament} % *GREEK
		% % --
			% \label{BELIEF-GREEK}
		
			% %
			% \subparagraph{\texorpdfstring{\gr{sample word}}{}}
			% %
				% \label{KATALLAGE}
			
				% \begin{compactdesc}
					% \item[BDAG , PDF ] \textbf{}
						% \begin{compactitem}
							% \item 
						% \end{compactitem}
					% \item[LSJ , PDF ] \textbf{}
						% \begin{compactitem}
							% \item 
						% \end{compactitem}
				% \end{compactdesc}
				
				% \begin{compactdesc}
					% \item[Some NT uses] \textbf{}
						% \begin{compactdesc}
							% \item[]
						% \end{compactdesc}
				% \end{compactdesc}
		
	% % -----
	% \subsubsection{Book of Mormon} % *BOM
	% % -----
		% \label{BELIEF-BOM}
	
		% \begin{tabular}{ll}
			% Total occurrences in the BOM & 
		% \end{tabular}
		
		% \begin{compactdesc}
			% \item[]
		% \end{compactdesc}
		
	% % -----
	% \subsubsection{Doctrine and Covenants} % *DC
	% % -----
		% \label{BELIEF-DC}
	
		% \begin{tabular}{ll}
			% Total occurrences in the DC & 
		% \end{tabular}
		
		% \begin{compactdesc}
			% \item[]
		% \end{compactdesc}
		
	% % -----
	% \subsubsection{Pearl of Great Price} % *PGP
	% % -----
		% \label{BELIEF-PGP}
	
		% \begin{tabular}{ll}
			% Total occurrences in the PGP & 
		% \end{tabular}
		
		% \begin{compactdesc}
			% \item[]
		% \end{compactdesc}
		
% % ----------------------
% \subsection{Key Phrases} % *KEYPHRASES *PHRASES
% % ----------------------
	% \label{BELIEF-PHRASES}

	% \begin{compactdesc}
		% \item[] \textbf{\#times} | list
			% % \begin{compactdesc}
				% % \item[Bible anchor verses] \phantom{.}
					% % \begin{compactdesc}
						% % \item[Romans 5:5] \gr{}
					% % \end{compactdesc}
				% % \item[Commentary] ---
			% % \end{compactdesc}
	% \end{compactdesc}
	
% % ----------------------
% \subsection{Bible Dictionaries} % *DICTIONARIES
% % ----------------------
	% \label{BELIEF-BD}

% % ----------------------
% \subsection{Restoration Usage} % *RESTORATION
% % ----------------------
	% \label{BELIEF-RESTORATION}

	% % -----
	% \subsubsection{Joseph Smith} % *JS
	% % -----
		% \label{BELIEF-JS}
	
		% \begin{compactdesc}
			% \item[{\hyperref[KEY]{Date/Source}}]
		% \end{compactdesc}
	
	% % -----
	% \subsubsection{Brigham Young} % *BY
	% % -----
		% \label{BELIEF-BY}
	
		% \begin{compactdesc}
			% \item[{\hyperref[KEY]{Date/Source}}]
		% \end{compactdesc}
	
% ----------------------
\subsection{Commentary and Studies} % *COMMENTARY *STUDIES
% ----------------------
	\label{BELIEF-COMMENTARY}

	% % -----
	% \subsubsection{Key Points}
	% % -----
		% \label{BELIEF-KEYS}
	
		% \begin{compactitem}
			% \item
		% \end{compactitem}
		
	% -----
	\subsubsection{Belief in belief and bearing testimony}
	% -----
		\label{BELIEF-belief}
		
		Lately (5/10/2021) I have been thinking about belief and belief in belief, which is something that Yudkowsky talks about in various places in his Less Wrong book. The point he makes, and others have made this too and I think this might have occurred to me independently a long time ago as well, is that in saying ``The Earth is round,'' what you are really saying is ``I believe the Earth is round.'' Each affirmation or denial you make should really carry with it the sign of that affirmation or denial. Going one step further, saying that ``I believe the Earth is round'' is really saying ``I believe I believe the Earth is round.''
		
		(There are a lot of other important points too---see parts B and H, with numbers 13, 17--19, and all of H in particular.)
		
		I think this is an important point in relation to religious belief in general, and I haven't thought through everything it means. But I have been considering its application to the practice of testimony-bearing. Virtually everyone who bears as testimony in an LDS setting actually says something like, ``I believe the church is true,'' or ``I believe that God exists.'' They actually include the ``I believe'' bit in what they say. Going by the rule of the previous paragraph, what they then really saying is ``I believe I believe that God exists.''
		
		This is helping me to understand what I dislike about the practice of testimony-bearing. Think about the difference in the internal processes required to say the words ``I believe that God exists'' versus ``God exists.'' Before you can affirm ``I believe that God exists,'' what you have to check is \textit{whether you actually believe that}. But that would be checking not whether something is true in the world or whether the world outside of you is a certain way, it would be checking to see whether you actually possess that belief. And if you look inside yourself and find that you do actually believe that God exists, you can then affirm, ``I believe that God exists.'' But the internal process shows that you are really affirming ``I believe I believe that God exists,'' because that's what you had to check.
		
		Compare affirming ``God exists.'' Before you can affirm that you have to check \textit{whether God actually exists}, which means you have to check something about the \textit{world} rather than about \textit{yourself}. It's like saying ``There is a computer on my desk.'' To affirm that, I don't have to check anything inside me, I have to go find out whether there really is a computer on my desk. If I discover that there is, I can then say the words. The difference in the processes is what you have to check before you make the affirmation.
		
		So in saying ``I believe that God exists,'' you are really just saying something about yourself: ``I believe I believe God exists.'' You are not making a claim about anything else in the world, it is only a claim about yourself. If other people wish to verify or investigate what you have said, they must check into \textit{you}, or your beliefs and your psychology, not into the world. And more than that, saying something like ``I believe that the church is true'' is really just a social signal---it communicates no useful content. It is the verbal equivalent of a thumbs-up or something like that.
		
		On the other hand, if I get up at the pulpit and say ``God exists,'' then I am in fact saying ``I believe that God exists.'' But I am affirming my belief \textit{about the world}, rather than \textit{about myself}. And so if someone wants to check whether I am right to believe that, they have to go check whether God exists, inasmuch as that is possible. They don't have to check into me, and it seems like that is the kind of testimony we want to be bearing anyway.